\documentclass{reportkit}

\setreportkitleftheader{REPORTKIT / ALGORITHM GRAPH VISUALIZATION TEST}
\setreportkitfooter{Algorithm graph/grid visualization acceptance test}
\setreportkitversion{v1.10.0}

\title{ReportKit -- graphstate and gridstate algorithm-visualization primitives}
\author{Test build}
\date{17 September 2026}

\begin{document}
\maketitle

\begin{execsummary}
This is an acceptance test for reportkit-algorithm-graph.sty's Phase 2 pair:
\texttt{graphstate} and \texttt{gridstate}
(docs/superpowers/specs/2026-09-16-reportkit-algorithm-visualization-primitives-spec.md,
sections 8-9). Both reuse the shared nine-state algorithm vocabulary
introduced in Phase 1 (see algorithm\_visuals\_acceptance\_test.tex) --
\texttt{graphstate} for traversal state over an explicit node/edge graph,
\texttt{gridstate} for traversal state over a 2D grid. Neither builds a
BFS-specific or DFS-specific primitive: a full traversal walkthrough is left
to the document author composing \texttt{graphstate} (or \texttt{gridstate})
with \texttt{algorithmtrace}, as the last section here demonstrates.
\end{execsummary}

\section{Graph traversal: BFS frontier (algorithm-interview scenario)}
A course-prerequisite graph mid-traversal: the start course has been
visited, its immediate neighbor is the node currently being processed, one
neighbor is queued on the frontier, and one course has not been reached yet.

\begin{diagram}[type=graph,width=\textwidth,caption={A breadth-first search frontier over a small prerequisite graph.},description={Four courses connected by prerequisite edges. Course A is visited, course B is the current node, course C is on the frontier, and course D is unseen.}]
\begin{graphstate}[columns=2]
  \graphnode[state=visited]{A}{Intro CS}
  \graphnode[state=current]{B}{Data Structures}
  \graphnode[state=frontier]{C}{Algorithms}
  \graphnode[state=unseen]{D}{Systems}
  \graphedge{A}{B}
  \graphedge{B}{C}
  \graphedge{B}{D}
\end{graphstate}
\end{diagram}

\section{Grid traversal: flood fill (algorithm-interview scenario)}
A flood fill expanding outward from a starting cell: visited cells have
already been colored in, the frontier is queued for the next expansion
step, and a wall cell blocks the fill from spreading further.

\begin{diagram}[type=grid,width=\textwidth,caption={A flood fill spreading from the top-left cell of a four by five grid.},description={A four row by five column grid. The top-left region is visited, its neighbors are on the frontier, one interior cell is currently being processed, and one cell is blocked by a wall.}]
\begin{gridstate}[rows=4,columns=5]
  \gridcell{1}{1}[state=visited]
  \gridcell{1}{2}[state=visited]
  \gridcell{2}{1}[state=visited]
  \gridcell{2}{2}[state=current]
  \gridcell{1}{3}[state=frontier]
  \gridcell{3}{2}[state=frontier]
  \gridcell{2}{3}[state=blocked]
  \gridcell{4}{5}[state=unseen]
\end{gridstate}
\end{diagram}

\section{Grid traversal: pipeline partition scan (data-engineering scenario)}
The same primitive applied outside interview algorithms: a partitioned
batch job scanning a data lake's row-group/column-chunk matrix for the
partitions that satisfy a filter predicate, with a corrupt chunk marked
blocked so it is skipped rather than silently dropped.

\begin{diagram}[type=grid,width=\textwidth,caption={A predicate scan across a partitioned dataset's row groups and column chunks.},description={A three row by six column grid representing row groups by column chunks. Scanned partitions are visited, the partition currently being read is current, the next two partitions are queued as frontier, and one chunk is blocked because it failed a checksum.}]
\begin{gridstate}[rows=3,columns=6]
  \gridcell{1}{1}[state=visited]
  \gridcell{1}{2}[state=visited]
  \gridcell{1}{3}[state=visited]
  \gridcell{2}{1}[state=visited]
  \gridcell{2}{2}[state=current]
  \gridcell{2}{3}[state=frontier]
  \gridcell{2}{4}[state=frontier]
  \gridcell{3}{2}[state=blocked]
\end{gridstate}
\end{diagram}

\section{Full node-state vocabulary}
All nine shared algorithm states applied to graph nodes, so state meaning
stays distinguishable through border weight and line style, not fill color
alone.

\begin{diagram}[type=graph,width=\textwidth,caption={Every shared algorithm state applied to one graph node each.},description={Nine graph nodes arranged on a grid, one per shared algorithm state: current, active, candidate, frontier, visited, resolved, discarded, blocked, and unseen.}]
\begin{graphstate}[columns=3]
  \graphnode[state=current]{n1}{current}
  \graphnode[state=active]{n2}{active}
  \graphnode[state=candidate]{n3}{candidate}
  \graphnode[state=frontier]{n4}{frontier}
  \graphnode[state=visited]{n5}{visited}
  \graphnode[state=resolved]{n6}{resolved}
  \graphnode[state=discarded]{n7}{discarded}
  \graphnode[state=blocked]{n8}{blocked}
  \graphnode[state=unseen]{n9}{unseen}
\end{graphstate}
\end{diagram}

\section{Composition: graphstate + algorithmtrace}
Per the spec's composition-before-specialization principle, a full BFS
walkthrough is expressed as ordered \texttt{graphstate} snapshots inside
\texttt{algorithmtrace}, not a dedicated BFS primitive.

\begin{diagram}[type=trace,width=\textwidth,caption={Three ordered snapshots of a breadth-first search expanding one level at a time.},description={Three ordered snapshots show a breadth-first search: the start node visited alone, then its two neighbors added to the frontier, then the first neighbor becoming current.}]
\begin{algorithmtrace}[columns=3]
  \snapshot{Start}{%
    \begin{graphstate}[columns=2,x spacing=2.2,y spacing=1.7]
      \graphnode[state=current]{a}{A}
      \graphnode[state=unseen]{b}{B}
      \graphnode[state=unseen]{c}{C}
      \graphedge{a}{b}
      \graphedge{a}{c}
    \end{graphstate}%
  }
  \snapshot{Discover}{%
    \begin{graphstate}[columns=2,x spacing=2.2,y spacing=1.7]
      \graphnode[state=visited]{a}{A}
      \graphnode[state=frontier]{b}{B}
      \graphnode[state=frontier]{c}{C}
      \graphedge{a}{b}
      \graphedge{a}{c}
    \end{graphstate}%
  }
  \snapshot{Advance}{%
    \begin{graphstate}[columns=2,x spacing=2.2,y spacing=1.7]
      \graphnode[state=visited]{a}{A}
      \graphnode[state=current]{b}{B}
      \graphnode[state=frontier]{c}{C}
      \graphedge{a}{b}
      \graphedge{a}{c}
    \end{graphstate}%
  }
\end{algorithmtrace}
\end{diagram}

\section{Dependency layout and ready queue (data-pipeline scenario)}
The algorithm-visuals fix sprint
(docs/superpowers/specs/2026-09-18-algorithm-visuals-fix-sprint-spec.md,
section 4) extends \texttt{dagstate} with \texttt{layout=dependency}: each
node's column becomes its longest-prerequisite distance (topological
level) instead of an arbitrary automatic-grid index, so a pipeline reads
left to right in execution order. \texttt{readyqueue} now renders an
actual FIFO queue box, labelled ``ready,'' rather than redrawing a node in
place -- and validates that every id it is given has indegree zero.

\begin{diagram}[type=dag,width=\textwidth,caption={A three-stage data pipeline mid-execution under Kahn's algorithm.},description={Three pipeline stages laid out by topological level: raw orders is resolved (already run), clean orders has no remaining unresolved dependency and is ready to run next, and daily sales still has one unresolved dependency. A ready queue box below the pipeline shows clean orders as the next node to execute.}]
\begin{dagstate}[layout=dependency]
  \dagnode[indegree=0,state=resolved]{raw}{Raw orders}
  \dagnode[indegree=0,state=frontier]{clean}{Clean orders}
  \dagnode[indegree=1,state=unseen]{sales}{Daily sales}
  \dependency{raw}{clean}
  \dependency{clean}{sales}
  \readyqueue{clean}
\end{dagstate}
\end{diagram}

\section{Composition: graphstate + queuestate (BFS frontier walkthrough)}
Per the spec's composition-before-specialization principle (spec section
4/C3), this sprint does not add a dedicated BFS primitive. A breadth-first
search work queue is instead expressed by placing an existing
\texttt{graphstate} beside an existing \texttt{queuestate} inside one
\texttt{diagram} -- proving the two independent primitives compose at
guide scale without a spacing or caption defect.

\begin{diagram}[type=graph,width=\textwidth,caption={A breadth-first search: the current node's unvisited neighbors enter the work queue.},description={A graph on the left shows a start node visited, its current neighbor being processed, one neighbor queued as frontier, and one neighbor unseen. A FIFO work queue on the right holds the two nodes discovered so far, front to rear.}]
\begin{graphstate}[columns=2]
  \graphnode[state=visited]{a}{Start}
  \graphnode[state=current]{b}{Current}
  \graphnode[state=frontier]{c}{Queued}
  \graphnode[state=unseen]{d}{Unseen}
  \graphedge{a}{b}
  \graphedge{b}{c}
  \graphedge{b}{d}
\end{graphstate}
\begin{scope}[xshift=8.6cm,yshift=-1.05cm]
  \begin{queuestate}
    \enqueue[state=visited]{Start}
    \enqueue[state=frontier]{Queued}
  \end{queuestate}
\end{scope}
\end{diagram}

\section{Composition: graphstate + stackstate (DFS frontier walkthrough)}
The equivalent depth-first search composition pairs \texttt{graphstate}
with \texttt{stackstate} instead of a dedicated DFS primitive.

\begin{diagram}[type=graph,width=\textwidth,caption={A depth-first search: the current node's unvisited neighbor is pushed onto the call stack.},description={A graph on the left shows a start node visited, its current neighbor being processed, one neighbor pushed as frontier, and one neighbor unseen. A LIFO call stack on the right holds the two nodes visited so far, with the most recently pushed node on top.}]
\begin{graphstate}[columns=2]
  \graphnode[state=visited]{e}{Start}
  \graphnode[state=current]{f}{Current}
  \graphnode[state=frontier]{g}{Pushed}
  \graphnode[state=unseen]{h}{Unseen}
  \graphedge{e}{f}
  \graphedge{f}{g}
  \graphedge{f}{h}
\end{graphstate}
\begin{scope}[xshift=8.6cm,yshift=-1.05cm]
  \begin{stackstate}
    \push[state=visited]{Start}
    \push[state=frontier]{Pushed}
  \end{stackstate}
\end{scope}
\end{diagram}

\section{Grayscale and semantic check}
\begin{principle}{State is not just color}
Every shared algorithm state carries a distinct line-style fragment in
addition to its draw/fill/text colors, so the full-vocabulary diagram above
should stay legible if printed in grayscale.
\end{principle}

\end{document}
